% Basic LaTeX2e Book Template — a plain `book`-class starter with the
% package/macro scaffolding an author actually reaches for: ntheorem-based
% examples, algorithm/algpseudocode, booktabs+dcolumn tables, natbib citations
% and microtype. Adapted from albertopessia/latex-book-template (public
% domain, Unlicense): dropped the upstream's dev-only packages (lipsum
% placeholder text, refcheck margin labels, lineno line numbers) since this is
% a document you ship, not one you are still drafting.
\documentclass[a4paper]{book}

\usepackage[a4paper]{geometry}

% Latin Modern font is an enhanced version of the default Computer Modern font
\usepackage{lmodern}
\usepackage[T1]{fontenc}
\usepackage[utf8]{inputenc}

% manuscript is written in British English (used for hyphenation)
\usepackage[british]{babel}

% shrink and stretch fonts for a good looking manuscript
\usepackage[activate={true, nocompatibility}, final, tracking=true, kerning=true, spacing=true, factor=1100, stretch=10, shrink=10]{microtype}

% bibliography management
\usepackage[round]{natbib}
\usepackage[nottoc, notlot, notlof]{tocbibind}

% improved version of amsmath
\usepackage{mathtools}
\usepackage{amssymb}

% customize theorem-like environments
\usepackage[amsmath, thref, thmmarks]{ntheorem}

% import figures
\usepackage{graphicx}

% typeset programs (programming code)
\usepackage{listings}

% algorithm environments
\usepackage{algpseudocode}
\usepackage{algorithm}

% easily customize the enumerate environment
\usepackage{enumitem}

% high quality tables
\usepackage{booktabs}
\usepackage{dcolumn}

% create hyperlinks for urls and references
\usepackage[colorlinks, citecolor=blue, urlcolor=blue, filecolor=blue]{hyperref}

\setcounter{tocdepth}{1}

\pagestyle{plain}

% add bold mathematical symbols to section titles
\makeatletter
\g@addto@macro\bfseries{\boldmath}
\makeatother

\renewcommand\thealgorithm{A.\arabic{algorithm}}

\algnewcommand{\Input}{\item[\textbf{Input:}]}
\algnewcommand{\Output}{\item[\textbf{Output:}]}

\newcommand{\clearemptydoublepage}{\clearpage{\pagestyle{empty}\cleardoublepage}}

\newcommand{\de}{\mathrm{d}}

\DeclareMathOperator*{\argmin}{arg\,min}
\DeclareMathOperator*{\argmax}{arg\,max}

\newcolumntype{d}[1]{D{.}{.}{#1}}

% Indented example environment
\theoremstyle{plain}
\theorembodyfont{\upshape}
\theoremseparator{.}
\theoremindent1\parindent
\newtheorem{example}{Example}[chapter]


\title{Notes on Algorithmic Number Theory}
\author{Alan Turing}
\date{}

\begin{document}

\frontmatter

\maketitle

\chapter{Preface}
These notes began as a set of lectures on algorithmic number theory,
written for students who already know basic abstract algebra but have not
yet seen how those ideas become running code. Each chapter pairs a short
piece of theory with the algorithm it justifies.

\chapter{Acknowledgements}
Thanks to the students of the original course for catching more errors in
early drafts than the author would like to admit.

\tableofcontents

\mainmatter

\chapter{Introduction}
\label{ch:introduction}

Modular arithmetic is the setting for most of what follows.
Definition~\ref{ex:modexp} recalls fast exponentiation, the workhorse
behind RSA and Diffie--Hellman alike \citep{cormen2009algorithms}.

\begin{example}[Fast modular exponentiation]
\label{ex:modexp}
To compute $a^n \bmod m$, write $n$ in binary and square-and-multiply:
repeated squaring reduces an $O(n)$-multiplication computation to
$O(\log n)$ multiplications modulo $m$.
\end{example}

Algorithm~\ref{alg:modexp} makes this precise.

\begin{algorithm}
\caption{Modular exponentiation by repeated squaring}
\label{alg:modexp}
\begin{algorithmic}
\Input{$a, n, m \in \mathbb{Z}$, $n \geq 0$}
\Output{$a^n \bmod m$}
\State $r \gets 1$
\State $b \gets a \bmod m$
\While{$n > 0$}
  \If{$n \bmod 2 = 1$}
    \State $r \gets (r \cdot b) \bmod m$
  \EndIf
  \State $b \gets (b \cdot b) \bmod m$
  \State $n \gets \lfloor n / 2 \rfloor$
\EndWhile
\State \Return $r$
\end{algorithmic}
\end{algorithm}

Table~\ref{tab:complexity} summarises the running time of the operations
this chapter introduces.

\begin{table}[htbp]
\centering
\caption{Running time of basic modular operations, $n$-bit inputs}
\label{tab:complexity}
\begin{tabular}{l l}
\toprule
Operation & Bit complexity \\
\midrule
Addition mod $m$       & $O(n)$ \\
Multiplication mod $m$ & $O(n^{1.58})$ \\
Exponentiation mod $m$ & $O(n^{2.58})$ \\
\bottomrule
\end{tabular}
\end{table}

\chapter{Primality}
\label{ch:primality}

\begin{example}[Fermat witness]
\label{ex:fermat}
If $\gcd(a, n) = 1$ and $a^{n-1} \not\equiv 1 \pmod{n}$, then $n$ is
composite and $a$ is called a \emph{Fermat witness} for $n$.
\end{example}

The Miller--Rabin test refines Example~\ref{ex:fermat} into a
probabilistic algorithm with one-sided error, described
in~\citet{cormen2009algorithms} and analysed asymptotically
in~\citet{knuth1997art}.

\backmatter

\bibliographystyle{abbrvnat}
\bibliography{lib/biblio}

\end{document}
